\documentclass[11pt,a4paper]{article}
\usepackage[utf8]{inputenc}
\usepackage[T1]{fontenc}
\usepackage{amsmath,amssymb}
\usepackage[margin=2.5cm]{geometry}
\usepackage{booktabs}
\usepackage{enumitem}
\usepackage[hidelinks]{hyperref}

\newcommand{\avg}[1]{\left\langle #1 \right\rangle}
\newcommand{\du}{\delta u}
\newcommand{\vect}[1]{\boldsymbol{#1}}

\title{The generalized Kolmogorov equation in particle-resolved channel flow:\\
status of the code and plan for a consistent treatment of the particles}
\author{GKE code, branch \texttt{dev\_xinju}}
\date{October 2026}

\begin{document}
\maketitle

\section{Purpose}

The branch \texttt{dev\_xinju} of the GKE code computes the generalized
Kolmogorov equation (GKE) of the particle-resolved channel-flow simulations of
X.~Jiang (the ``xinju dataset''), with the averages restricted to pairs of
fluid points. This document summarizes what is computed today, why the
resulting budget is not yet closed, and how the action of the particles on the
fluid can be included consistently. It is meant as the working plan for the
next development steps.

\section{The flow solver and the stored data}
\label{sec:solver}

The simulations follow the numerical method described in Appendix~A of
Jiang, Xu \& Zhao (2024). The fluid obeys the incompressible Navier--Stokes
equations in non-conservative form, extended to the whole computational box by
the immersed-boundary (IB) force $\tilde{\vect f}$, which has the units of an
acceleration:
\begin{equation}
  \nabla\cdot\tilde{\vect u}=0,\qquad
  \frac{\partial\tilde{\vect u}}{\partial t}+\tilde{\vect u}\cdot\nabla\tilde{\vect u}
  =-\nabla\tilde p+\nu\nabla^2\tilde{\vect u}+\tilde{\vect f}.
  \label{eq:ns}
\end{equation}
Tildes denote total quantities, and the density is absorbed into $\tilde p$. The
discretization uses second-order central finite differences on a uniform
staggered grid ($\Delta x=\Delta y=\Delta z=\Delta h$) and a fractional-step
method. A first prediction $\vect u^*$ is computed with Crank--Nicolson for the
viscous term and without the IB force. The IB force is then added,
$\vect u^{**}=\vect u^*+\Delta t\,\vect f^{n+1/2}$, followed by the projection
that updates pressure and velocity. The IB force is computed by direct forcing
on Lagrangian markers that are retracted by $0.3\Delta h$ inside the particle
surface, with two multi-forcing iterations. It is spread to the Eulerian grid
with the three-cell regularized delta function of Roma, Peskin \& Berger
(1999), whose support is $3\Delta h$ centred on the markers.

The following properties of the stored snapshots were verified on the
snapshot $n=20000$ ($480\times160\times240$ cells, $\Delta h=0.0125$, solid
volume fraction about $0.30$):
\begin{itemize}[itemsep=1pt]
\item each MPI-rank file stores the 11 variables $u,v,w,f^\mathrm{turb}_{x,y,z},
  f^\mathrm{IB}_{x,y,z},p,\phi$ as consecutive blocks (Fortran order);
  $f^\mathrm{turb}\equiv0$;
\item the discrete divergence of the staggered velocity vanishes to round-off
  ($\approx10^{-14}$) everywhere, inside the particles too, as expected
  from~\eqref{eq:ns};
\item the level set $\phi$ is negative in the solid and positive in the fluid,
  but it is not a distance: $|\nabla\phi|\approx10$ near the surface, and it
  saturates at $0.3$ in the fluid;
\item $\tilde{\vect f}$ is concentrated in a layer about two cells wide on
  the fluid side and three cells wide on the solid side of the surface, as the
  marker retraction and the kernel support imply. Further away it is small:
  rms between $10^{-3}$ and $10^{-2}$, against about $1.5$ in the layer.
\end{itemize}
The conversion script \texttt{convert\_h5.py} distributed with the data
reshapes the rank files with the variable as the fastest index, so its HDF5
output is scrambled. The GKE code reads the rank files directly and is not
affected.

\section{What the code computes today}
\label{sec:current}

Fluctuations are defined with respect to the mean velocity $(U,V,W)(y)$ and
pressure $P(y)$ of the fluid phase, i.e.\ averaged over the fluid points of
each wall-parallel plane and over the snapshots. For two points
$\vect x_1=\vect X-\vect r/2$ and $\vect x_2=\vect X+\vect r/2$ we write
$\du_i=u_i(\vect x_2)-u_i(\vect x_1)$, $u^*_i=[u_i(\vect x_1)+u_i(\vect x_2)]/2$,
$q=\du_i\du_i$, and analogously $\delta U_i$, $V^*$ and $(\cdot)^*$ for the
mean quantities. Every average $\avg{\cdot}$ is taken over pairs of fluid
points only and renormalized by the number of such pairs. With
$\varepsilon=\nu\avg{\partial_k u_i\partial_k u_i}$ the pseudo-dissipation,
the code evaluates
\begin{align}
  \Phi_{r_k}&=\avg{\du_k q}+\delta U_k\avg{q}-2\nu\frac{\partial\avg{q}}{\partial r_k},
  \label{eq:phir}\\
  \Phi_Y&=\avg{v^*q}+V^*\avg{q}+2\avg{\delta p\,\du_2}-\frac{\nu}{2}\frac{\partial\avg{q}}{\partial Y},
  \label{eq:phic}\\
  \xi&=-2\avg{\du_i\du_2}\Big(\frac{dU_i}{dy}\Big)^*
       -2\avg{\du_i v^*}\,\delta\Big(\frac{dU_i}{dy}\Big)-4\varepsilon^*,
  \label{eq:xi}
\end{align}
which make up the single-phase GKE
$\partial_t\avg{q}+\nabla_{\vect r}\cdot\vect\Phi_{\vect r}+\partial_Y\Phi_Y=\xi$
(Hill 2002; Marati, Casciola \& Piva 2004; Gatti et al.\ 2019). In
\texttt{gke.bin} they are stored as \texttt{phiR}, \texttt{phiC},
\texttt{scaleENER}$=\avg{q}$ and \texttt{scalePROD}$=\xi$. The pair sums are
evaluated exactly with FFT cross-correlations of mask-weighted fields
(\texttt{step2\_gke}). The direct pair accumulation \texttt{step2\_gke\_direct}
serves as reference, and the two agree to round-off.

\paragraph{Why this budget is not closed.} Equations
\eqref{eq:phir}--\eqref{eq:xi} are the GKE of a single-phase flow. Restricted
to fluid pairs, they miss:
\begin{enumerate}[label=(\alph*),itemsep=1pt]
\item the IB force, i.e.\ the source $2\avg{\du_i\,\delta f_i}$, which is
  active in the forcing layer, partly inside the fluid;
\item the interface terms. The pair indicator does not commute with the
  derivatives in $\vect r$, $Y$ and $t$, which produces fluxes through the
  particle surfaces (section~\ref{sec:fluid});
\item the effect of the normalization. The fraction of fluid pairs depends on
  $(\vect r,Y)$ (the solid fraction ranges from about $1\%$ at the walls to
  $50\%$ at the centre), so derivatives of renormalized averages contain
  gradients of that fraction;
\item the conditional-mean terms. Over fluid pairs, $\avg{\du_i}$ and
  $\avg{\delta p}$ do not vanish, so terms such as
  $-2\avg{\delta p}\,\delta(dV/dy)$ and the products of $\avg{\du_i}$ with
  the mean-momentum terms survive. These are neglected at present.
\end{enumerate}
Items (b)--(d) disappear for unmasked averages, and item (a) is then the only
missing term. This motivates the two-stage plan below.

\section{Formulation A: whole-domain GKE with the IB force}
\label{sec:whole}

Since~\eqref{eq:ns} holds at every grid point, the GKE of the whole-domain
field is exact once the IB force is included. With averages over all points
(fluid and solid), Reynolds decomposition
$\tilde{\vect u}=\vect U(y)+\vect u$, $\tilde p=P(y)+p$,
$\tilde{\vect f}=\vect F(y)+\vect f$, and the usual steps, one obtains
\begin{equation}
  \frac{\partial\avg{q}}{\partial t}+\nabla_{\vect r}\cdot\vect\Phi_{\vect r}
  +\frac{\partial\Phi_Y}{\partial Y}=\xi+\Pi_f,\qquad
  \Pi_f=2\avg{\du_i\,\delta f_i},
  \label{eq:gkef}
\end{equation}
with the fluxes and the source of~\eqref{eq:phir}--\eqref{eq:xi}. Continuity
and the walls imply $V\equiv0$ for whole-domain averages; uniform forces (the
driving pressure gradient, the buoyancy-corrected weight) cancel in the
increments. The single-point budgets gain
$\avg{u_if_j}+\avg{u_jf_i}$, and the mean momentum balance reads
\begin{equation}
  \frac{d}{dy}\Big(\nu\frac{dU}{dy}-\avg{uv}\Big)+\Pi+F_x=0,
  \label{eq:mean}
\end{equation}
where $\Pi$ is the (uniform) driving pressure gradient and $F_x$ the mean IB
force, i.e.\ the mean particle drag per unit mass.

\paragraph{Interpretation.} $\Pi_f(\vect r,Y)$ distributes the energy
exchange between particles and fluid over scales and positions. At large
separations it tends to the sum of the single-point exchanges at the two
points, $2\avg{u_if_i}_1+2\avg{u_if_i}_2$. The price is that the averages
include the solid interior, where $\tilde{\vect u}$ is (approximately) the
rigid-body motion, so the statistics describe the fluid--particle mixture
rather than the fluid alone.

\paragraph{Changes to the code.}
\begin{itemize}[itemsep=1pt]
\item \texttt{dataset.cpl}: read $f^\mathrm{IB}$ (variables 7--9), staggered
  like the velocity, and interpolate it to the cell centres.
\item \texttt{gke.in}: option to switch the masking off (and, for
  section~\ref{sec:fluid}, to choose the mask of each point and the
  definition of the mean).
\item Step 1: mean force $\vect F(y)$, correlations $\avg{u_if_j}$ (new field
  of \texttt{BALANCE}).
\item Step 2: one more pair sum $\sum\du_i\,\delta f_i$. With $\vect a$ the
  velocity fluctuation it reads $P(1,\vect a\cdot\vect f)-\sum_iP(a_i,f_i)$
  in the notation of \texttt{step2\_gke.cpl}, i.e.\ four more spectra (about
  $0.6$\,GiB for the xinju grid). New fields of \texttt{GKETERMS}: $\Pi_f$,
  and the fraction of accumulated pairs, which is needed in
  section~\ref{sec:fluid}.
\end{itemize}

\section{Formulation B: fluid-phase GKE with interface terms}
\label{sec:fluid}

Let $m(\vect x,t)$ be the fluid indicator and
$\chi=m(\vect x_1,t)\,m(\vect x_2,t)$ the pair indicator. Write the exact pair
equation in pair space $(\vect X,\vect r)$ as
$\partial_tq+\nabla\cdot\vect J=s$. Here $\vect J$ collects the advective,
pressure and viscous fluxes and $s$ the production, dissipation and forcing.
Multiplying by $\chi$ and averaging gives the budget of the phase-weighted
quantity $\avg{\chi q}$ (averages over \emph{all} pairs):
\begin{equation}
  \partial_t\avg{\chi q}+\nabla\cdot\avg{\chi\vect J}=\avg{\chi s}+I,\qquad
  I=\avg{q\,\partial_t\chi+\vect J\cdot\nabla\chi}.
  \label{eq:phase}
\end{equation}
The particle surfaces move with the rigid-body velocity $\vect u_\Gamma$,
$\partial_tm+\vect u_\Gamma\cdot\nabla m=0$, and no-slip gives
$\tilde{\vect u}=\vect u_\Gamma$ there. The temporal and advective parts of $I$
therefore cancel, and $I$ reduces to the pressure and viscous pair fluxes
through the particle surfaces. These are weighted by the other point of the
pair being in the fluid, since $\nabla m=\vect n\,\delta_\Gamma$, with
$\vect n$ the normal pointing into the fluid.

In the IB discretization the surface stresses are not sharp. They are
represented by $\tilde{\vect f}$ in the forcing layer, which straddles the
surface. Hence $I$ and the forcing contribution $\avg{\chi\,2\du_i\delta f_i}$
cannot be separated meaningfully and should be reported together, as the
\emph{fluid--particle exchange} of the fluid-phase budget.

\paragraph{Computation through a decomposition into pair classes.} No surface
integrals are needed. With $m_\mathrm{f}=m$ and $m_\mathrm{s}=1-m$, the four
pair classes
$\chi_{ab}=m_a(\vect x_1)m_b(\vect x_2)$, $a,b\in\{\mathrm{f},\mathrm{s}\}$,
add up to one. Their phase-weighted budgets~\eqref{eq:phase} therefore add up
to the closed whole-domain budget~\eqref{eq:gkef}, and their interface terms
$I_{ab}$ are exchanges between classes that sum to zero. Every term of every
class is computed with the existing FFT machinery, by choosing the mask of
each of the two points. The fluid--particle exchange of the fluid--fluid
class follows as the residual of its budget. This assumes statistical
stationarity, i.e.\ enough snapshots for $\partial_t\avg{\chi q}$ to vanish.

\paragraph{Choices to be made.}
\begin{itemize}[itemsep=1pt]
\item \emph{Reference mean.} The classes add up only if all of them use the
  same fluctuations. The whole-domain means ($V\equiv0$) are the natural
  choice. Fluid-phase means can still be used for presentation, at the price
  of the conditional-mean terms of item~(d).
\item \emph{Normalization.} The budget is formulated for $\avg{\chi q}$.
  Renormalized averages $\avg{q}_\chi=\avg{\chi q}/\avg{\chi}$ are obtained
  by dividing only at the end. Their budget follows from the quotient rule
  and contains $\nabla\avg{\chi}$. The pair fraction $\avg{\chi}(\vect r,Y)$
  must therefore be stored.
\end{itemize}

\section{Numerical consistency and validation}
\label{sec:validation}

Closure tests are the core of the validation. Expected residuals stem from:
\begin{itemize}[itemsep=1pt]
\item \emph{time levels}: the fractional step stores fields at different time
  levels (e.g.\ $\vect u^{n+1}$, $p^{n+1/2}$, $\vect f^{n+1/2}$), an
  $O(\Delta t)$ inconsistency, small at the CFL numbers used;
\item \emph{spatial discretization}: the GKE code evaluates derivatives with
  second-order finite differences of the cell-centred fields. The diagonal
  derivatives use the solver's own differences across the cell faces, so the
  discrete divergence vanishes, but discrete product rules hold only up to
  $O(\Delta h^2)$. Near the surfaces the fields are only continuous, so
  larger local errors are expected there;
\item \emph{statistics}: the time derivative vanishes only on average.
\end{itemize}
The suggested sequence is:
\begin{enumerate}[itemsep=1pt]
\item \emph{Mean momentum}~\eqref{eq:mean}: the driving gradient $\Pi$
  obtained as residual must be independent of $y$. This tests the sign,
  units and staggering of $\tilde{\vect f}$ with almost no effort.
\item \emph{Single-point energy budget} (whole domain): the residual of the
  budget of $\avg{u_iu_i}$, including $2\avg{u_if_i}$, must be small compared
  with its largest terms.
\item \emph{GKE~\eqref{eq:gkef}}: a post-processing script evaluating
  $\nabla_{\vect r}\cdot\vect\Phi_{\vect r}$ on the undersampled separations
  and $\partial_Y\Phi_Y$, and the residual map in $(\vect r,Y)$.
\item \emph{Pair classes}: the four class budgets must add up to the
  whole-domain one to round-off. The fluid--fluid residual gives the
  fluid--particle exchange.
\end{enumerate}

\section{Implementation steps}

\begin{enumerate}[itemsep=1pt]
\item Read $f^\mathrm{IB}$; masking switch in \texttt{gke.in}; mean
  momentum check (validation step 1).
\item Force terms in steps 1 and 2; new fields of \texttt{BALANCE} and
  \texttt{GKETERMS} (pair fraction included); single-point closure check.
\item Post-processing script for the GKE closure; validation on the synthetic
  test case (with a synthetic force field) and on the xinju dataset.
\item Per-point masks (fluid, solid, all) and the whole-domain reference mean
  as options; class decomposition and fluid--fluid exchange.
\item Falling particles. With walls normal to gravity nothing changes
  structurally, but the fluid-phase mean $V(y)$ is non-zero (already
  supported). If the flow is periodic in all directions, as in Jiang et al.\
  (2024), the statistics are homogeneous: the GKE depends on $\vect r$ only,
  and the wall-normal grid, the mirror symmetry and step~3 have to be replaced
  by a fully periodic variant.
\end{enumerate}

\section{Open questions for the data provider}

\begin{enumerate}[itemsep=1pt]
\item Time levels of the stored $\vect u$, $p$ and $\vect f^\mathrm{IB}$.
  Is $\vect f^\mathrm{IB}$ the total force of the two multi-forcing
  iterations?
\item Is the pressure gradient that keeps the flow rate fixed contained in
  $p$, or applied as a separate uniform forcing? Is its value per snapshot
  available?
\item Definition of the level set (it is not a distance, $|\nabla\phi|\approx10$).
  Does $\phi=0$ coincide with the particle surface or with the retracted
  markers?
\item Are the particle positions, velocities and angular velocities available
  per snapshot (useful for the rigid-body checks and the solid-phase
  statistics)? What are the particle size, the density ratio and the
  direction of gravity in the channel cases?
\end{enumerate}

\section*{References}
\begin{list}{}{\leftmargin=1.5em\itemindent=-1.5em\itemsep=1pt}
\item Gatti, D. et al. 2019
  An efficient numerical method for the generalized Kolmogorov equation.
  \emph{J. Turbul.} \textbf{20}, 457--480.
\item Hill, R.J. 2002 Exact second-order structure-function relationships.
  \emph{J. Fluid Mech.} \textbf{468}, 317--326.
\item Jiang, X., Xu, C. \& Zhao, L. 2024 Prolate spheroids settling in a
  quiescent fluid: clustering, microstructures and collisions. \emph{J. Fluid
  Mech.} \textbf{1000}, A49.
\item Marati, N., Casciola, C.M. \& Piva, R. 2004 Energy cascade and spatial
  fluxes in wall turbulence. \emph{J. Fluid Mech.} \textbf{521}, 191--215.
\item Roma, A.M., Peskin, C.S. \& Berger, M.J. 1999 An adaptive version of the
  immersed boundary method. \emph{J. Comput. Phys.} \textbf{153}, 509--534.
\end{list}

\end{document}
